\documentclass{article}
\usepackage{amsmath, amssymb, amsthm}

\begin{document}

\noindent\textbf{Claim.} The robot can never reach $(1,1)$.

\begin{proof}
Define $f(x,y) = 4x + y$. We claim that
\[
f(x,y) \equiv 0 \pmod{7}
\]
is an invariant of the robot's motion.

\medskip\noindent\textbf{Base case:} At the starting position $(0,0)$,
\[
f(0,0) = 0 \equiv 0 \pmod{7}.
\]

\medskip\noindent\textbf{Inductive step:} Suppose the robot is at a position $(x,y)$ such that
\[
f(x,y) \equiv 0 \pmod{7}.
\]
We show that after one move, the invariant is preserved. Consider each possible move:

\begin{align*}
(+2,-1): \quad f(x+2,y-1) &= 4(x+2) + (y-1) \\
&= 4x + y + 8 - 1 \\
&= f(x,y) + 7 \\
&\equiv 0 \pmod{7}.
\end{align*}

\begin{align*}
(-2,+1): \quad f(x-2,y+1) &= 4(x-2) + (y+1) \\
&= 4x + y - 8 + 1 \\
&= f(x,y) - 7 \\
&\equiv 0 \pmod{7}.
\end{align*}

\begin{align*}
(+1,+3): \quad f(x+1,y+3) &= 4(x+1) + (y+3) \\
&= 4x + y + 4 + 3 \\
&= f(x,y) + 7 \\
&\equiv 0 \pmod{7}.
\end{align*}

\begin{align*}
(-1,-3): \quad f(x-1,y-3) &= 4(x-1) + (y-3) \\
&= 4x + y - 4 - 3 \\
&= f(x,y) - 7 \\
&\equiv 0 \pmod{7}.
\end{align*}

\noindent Thus, in all cases, the invariant is preserved. 
By induction on the number of moves, every reachable position $(x,y)$ satisfies
\[
f(x,y) \equiv 0 \pmod{7}.
\]
However,
\[
f(1,1) = 4 + 1 = 5 \not\equiv 0 \pmod{7}.
\]
Therefore, the position $(1,1)$ is unreachable.

\end{proof}

\end{document}
